\section{Introduction}

Émile Zola, l'un des écrivains les plus célèbres du XIXe siècle, est connu pour ses romans naturalistes qui explorent les aspects sociaux, économiques et psychologiques de la vie humaine. 

Dans cette analyse, nous allons examiner les thèmes récurrents dans les œuvres de Zola et déterminer s'ils apparaissent dans plusieurs de ses romans. Des méthodes d'analyse textuelle seront utilisées pour identifier les thèmes principaux à travers les différents romans.

La question centrale de cette étude est donc la suivante : dans quelle mesure des méthodes non supervisées comme le NMF et K-means permettent-elles de faire apparaître des structures thématiques récurrentes dans l'œuvre romanesque de Zola ?

L'objectif est de vérifier si les regroupements produits automatiquement correspondent à des motifs interprétables dans le cadre d'une analyse littéraire.

Ce travail est le premier d'une série d'analyses portant sur un corpus de textes littéraires issus de différents auteurs naturalistes et non naturalistes.

Deux méthodes d'analyse ont été utilisées pour identifier les thèmes dans les romans de Zola : 
\vspace{0.5cm}
\begin{itemize}
\item le NMF (Non-negative Matrix Factorization)\footnote{N. Gillis, \textit{Nonnegative Matrix Factorization}, Society for Industrial and Applied Mathematics, coll. « Data Science », 2020.} : une technique de réduction de dimensionnalité qui permet d'extraire des thèmes à partir d'un corpus de textes. 

\item K-means : une méthode de clustering qui permet de regrouper les thèmes récurrents en fonction de leur similarité thématique.
\end{itemize}
\vspace{0.5cm}
Ces méthodes seront expliquées en détail dans les sections suivantes, où nous présenterons les résultats de l'analyse thématique des romans de Zola.

